\chapter{Un éditeur de texte, puis un éditeur de code}

Un éditeur de texte est le premier programme où l'utilisateur \emph{modifie} des
données au lieu de les regarder. Tout change : il faut un curseur, une sélection,
un défilement, une annulation --- et il faut décider comment le texte est rangé
en mémoire, décision qui gouvernera tout le reste.

Le dépôt contient un éditeur mature, celui de NKCode : \textbf{5\,012 lignes} pour
le seul \texttt{NkCodeEditor.h}. Ce chapitre en extrait les quatre décisions qui
comptent, puis vous fait construire le vôtre --- d'abord texte, puis code.

\section{Décision 1 : comment ranger le texte}

\begin{nkcode}{Applications/NKCode/src/NKCode/Editor/NkCodeEditor.h}
// Une ligne = caracteres bruts (PAS de '\0').
using NkCodeLine = NkVector<char>;

struct NkCodeDoc {
        NkVector<NkCodeLine> lines;
        int32 curLine = 0, curCol = 0;   // curseur
        int32 selLine = 0, selCol = 0;   // ancre de selection
        // ...
};
\end{nkcode}

Un document est un \emph{vecteur de lignes}, chaque ligne un \emph{vecteur de
caractères}.

\begin{nkretenir}
\textbf{Il y a trois façons courantes de ranger un texte éditable}, et le choix se
juge sur ce qu'il rend cher :

\begin{center}
\begin{tabular}{@{}lll@{}}
\toprule
\textbf{Représentation} & \textbf{Bon marché} & \textbf{Cher} \\
\midrule
une seule chaîne & lire tout d'un coup & insérer au milieu \\
vecteur de lignes & éditer une ligne, dessiner & insérer une ligne au début \\
tampon à trou (\emph{gap}) & frappe au clavier & sauts lointains \\
\bottomrule
\end{tabular}
\end{center}

NKCode prend le \textbf{vecteur de lignes}, et c'est cohérent avec ce qu'il fait :
un éditeur passe son temps à \emph{dessiner des lignes} et à \emph{éditer celle
du curseur}. L'opération chère --- décaler un million de lignes --- n'arrive
jamais.

\textbf{Choisissez d'après l'opération la plus fréquente, pas d'après la plus
élégante.}
\end{nkretenir}

\begin{nkpiege}
\textbf{« PAS de \texttt{'\textbackslash 0'} »} : le commentaire d'origine le dit
en majuscules, et c'est un piège réel. Une ligne est une plage
$[\texttt{begin}, \texttt{end})$, pas une chaîne C. Passer \texttt{line.Data()} à
une fonction qui attend une chaîne terminée lirait au-delà de la ligne, jusqu'au
premier zéro rencontré dans le tas --- donc parfois rien d'anormal, et parfois un
plantage. Le pire des deux mondes.
\end{nkpiege}

\section{Décision 2 : où vit l'état d'édition}

Le curseur, la sélection et le défilement sont \emph{dans} \texttt{NkCodeDoc},
avec le texte. L'en-tête explique pourquoi :

\begin{nkcode}{Applications/NKCode/src/NKCode/Editor/NkCodeEditor.h}
// etat curseur/selection/scroll EMBARQUE dans le document -> survit au
// changement d'onglet ET au re-dock (etat hors-frame).
\end{nkcode}

\begin{nkretenir}
\textbf{L'interface est en mode immédiat : elle ne garde rien d'une image à
l'autre.} Un état rangé dans le widget disparaîtrait au premier changement
d'onglet --- vous reviendriez sur votre fichier, curseur au début, défilement en
haut.

Donc : \textbf{ce qui doit survivre à la disparition du widget vit dans le
document, pas dans le widget}. C'est la règle générale du mode immédiat, et
l'éditeur en est le cas le plus visible.
\end{nkretenir}

\section{Décision 3 : l'annulation}

\begin{nkcode}{Applications/NKCode/src/NKCode/Editor/NkCodeEditor.h}
struct UndoState { /* ... */ };
\end{nkcode}

\begin{nkretenir}
\textbf{Deux stratégies, et la plus simple est souvent la bonne.}

\emph{Instantané} : on garde une copie du document avant chaque groupe de
modifications. Simple, sûr, et coûteux en mémoire --- mais un fichier source fait
quelques dizaines de kilo-octets, donc cent états tiennent dans quelques
méga-octets.

\emph{Journal d'opérations} : on garde « inséré \texttt{x} en (ligne, colonne) »
et on sait le défaire. Économe, et bien plus difficile à rendre juste : chaque
opération a besoin de son inverse exact, et un seul inverse faux corrompt tout
l'historique sans prévenir.

\textbf{Commencez par l'instantané.} Ne passez au journal que si une mesure ---
pas une intuition --- montre que la mémoire pose problème.
\end{nkretenir}

\begin{nkpiege}
\textbf{Un groupe d'annulation n'est pas une frappe.} Si chaque caractère tapé
crée un état, l'utilisateur devra faire trente annulations pour effacer un mot.

Le regroupement se fait sur un \emph{critère} : une pause de frappe, un changement
de ligne, un changement de nature d'opération (saisie puis effacement). Ce critère
est un choix de conception, pas un détail --- et c'est celui que les utilisateurs
remarquent en premier.
\end{nkpiege}

\section{Décision 4 : ne dessiner que ce qui se voit}

Un fichier de cent mille lignes s'affiche dans une fenêtre qui en montre
cinquante. Dessiner les cent mille serait deux mille fois trop de travail --- et
la fluidité s'effondrerait sur le premier gros fichier.

\begin{nkretenir}
\textbf{On calcule la première et la dernière ligne visibles, et on ne dessine
que celles-là.} Le coût du dessin devient indépendant de la taille du fichier :
cinquante lignes, qu'il y en ait cent ou cent mille.

L'en-tête de NKCode le dit d'ailleurs pour la largeur : « largeur H approximée
sur les lignes visibles ». Même la largeur de défilement horizontal n'est calculée
que sur ce qui est à l'écran --- une approximation \emph{assumée et écrite}, parce
que mesurer cent mille lignes à chaque image coûterait plus que le bénéfice.
\end{nkretenir}

\section{De l'éditeur de texte à l'éditeur de code}

La différence tient en un mot : \textbf{la coloration}. Et la façon dont NKCode
l'aborde est la vraie leçon de cette seconde moitié.

\begin{nkcode}{Applications/NKCode/src/NKCode/Editor/NkSyntax.h}
struct NkLangSpec {
        NkString name;
        NkVector<NkString> exts;         ///< ".css", ".scss"...
        NkString lineComment;            ///< "//", "#", "--", ";" ("" = aucun)
        NkString blockOpen, blockClose;  ///< "/*","*/" ou "<!--","-->"
        NkString strChars;               ///< delimiteurs de chaines
        NkVector<NkString> keywords;     ///< tries
        NkVector<NkString> types;        ///< tries
};
\end{nkcode}

\begin{nkretenir}
\textbf{Un langage n'est pas du code : c'est une donnée.} Ajouter la coloration du
Lua, du Rust ou du CSS ne demande pas d'écrire un colorateur --- seulement de
remplir un \texttt{NkLangSpec}.

C'est exactement le principe que le cours répète depuis le début : \emph{décrire
plutôt que programmer}. Un colorateur écrit langage par langage aurait donné dix
fonctions qui divergent ; une description en donne dix \emph{données} qui ne
peuvent pas diverger, puisqu'un seul code les lit.
\end{nkretenir}

\begin{nkcode}{Applications/NKCode/src/NKCode/Editor/NkSyntax.h}
// Tokenise UNE ligne et emet des plages colorees [begin,end) via un callback
// Sans allocation : gap-filling (les zones non colorees -> couleur texte)
\end{nkcode}

\begin{nkretenir}
\textbf{Deux décisions de performance dans deux lignes de commentaire.}

\emph{Une ligne à la fois} : la coloration se fait au moment de dessiner, sur les
lignes visibles seulement. Rien n'est stocké, donc rien ne peut être périmé après
une modification.

\emph{Sans allocation, par remplissage de trous} : au lieu de produire une liste
de toutes les plages --- y compris celles qui gardent la couleur par défaut --- le
tokeniseur n'émet que les plages \emph{colorées}, et ce qui reste entre elles
prend la couleur du texte. Moins d'objets, aucune allocation par image.
\end{nkretenir}

\begin{nkpiege}
\textbf{Un commentaire de bloc ne tient pas dans une ligne.} \texttt{/*} ouvert
ligne 10 et fermé ligne 40 : les trente lignes du milieu sont dans le commentaire,
et rien dans leur texte ne le dit.

Un tokeniseur strictement par ligne ne peut pas le savoir : il lui faut un
\textbf{état d'entrée de ligne} --- « suis-je déjà dans un commentaire ? » ---
calculé depuis le début du fichier, ou depuis un point de reprise connu.

C'est la seule vraie difficulté de la coloration syntaxique, et c'est celle qu'on
découvre en dernier, quand un fichier entier devient vert.
\end{nkpiege}

\begin{nkbilan}
Un éditeur repose sur quatre décisions : la représentation du texte (choisie
d'après l'opération \emph{la plus fréquente} --- ici un vecteur de lignes), le
lieu de l'état d'édition (dans le document, jamais dans le widget, parce que le
mode immédiat ne garde rien), l'annulation (l'instantané avant le journal, et le
regroupement est ce que l'utilisateur remarque), et le dessin restreint aux lignes
visibles. Passer au code n'ajoute qu'une chose : la coloration --- et le bon
réflexe est de \emph{décrire} chaque langage par une donnée plutôt que de le
programmer.
\end{nkbilan}

\begin{nkquestions}
\begin{enumerate}
\item Citez trois représentations d'un texte éditable et l'opération que chacune
      rend chère.
\item Pourquoi l'état du curseur ne peut-il pas vivre dans le widget ?
\item Pourquoi une ligne de NKCode ne se termine-t-elle pas par
      \texttt{'\textbackslash 0'} ?
\item Quel est l'inconvénient d'un groupe d'annulation par caractère tapé ?
\item Pourquoi un tokeniseur par ligne a-t-il besoin d'un état d'entrée de ligne ?
\end{enumerate}
\end{nkquestions}

\begin{nkpratique}
Construisez votre éditeur de texte, en deux étapes.

\textbf{Étape A --- le texte.} Un \texttt{NkCanvasGuiApp} qui ouvre un fichier, le
range en vecteur de lignes, affiche les lignes visibles avec leurs numéros, place
un curseur, permet les flèches, \texttt{Home}, \texttt{End}, la saisie,
\texttt{Entrée}, \texttt{Retour arrière} et \texttt{Suppr}, défile à la molette,
et enregistre.

\textbf{Étape B --- le code.} Ajoutez un \texttt{NkLangSpec} pour un langage de
votre choix et colorez : mots-clés, types, chaînes, nombres, commentaires de
ligne. Puis ajoutez un \emph{second} langage \textbf{sans écrire une seule ligne
de code de coloration} --- seulement une donnée. Si vous n'y arrivez pas, votre
étape B n'était pas assez « décrite ».
\end{nkpratique}

\begin{nkdemo}
Prouvez que votre éditeur ne dessine que les lignes visibles.

Ajoutez un compteur de lignes réellement dessinées, affiché à l'écran. Ouvrez un
fichier de dix lignes, puis un fichier de cent mille lignes générées. Le compteur
doit être \emph{le même} dans les deux cas.

Mesurez ensuite le temps d'une image avec \texttt{NkChrono}, sur les deux
fichiers, dix fois chacun.

\textbf{Ce que la démonstration doit établir} : que l'écart entre les deux temps
reste sous le bruit de mesure --- et il vous faut donc \emph{d'abord} mesurer ce
bruit, en comparant deux séries sur le \emph{même} fichier. Sans ce témoin, vous
ne pouvez pas dire si un écart de 3\,\% est un effet ou du hasard.
\end{nkdemo}

\begin{nklogique}
Un utilisateur tape \texttt{Bonjour}, appuie sur \texttt{Entrée}, tape
\texttt{monde}, puis fait trois annulations.

\begin{enumerate}
\item Décrivez l'état obtenu selon trois politiques de regroupement : par
      caractère, par mot, par ligne.
\item Laquelle vous paraît la meilleure ? Défendez-la, puis donnez le cas où elle
      surprend l'utilisateur.
\item Le document est ensuite \emph{rechargé depuis le disque} par une action
      extérieure. L'historique d'annulation doit-il survivre ? Argumentez dans les
      deux sens, puis tranchez et dites ce que votre choix rend impossible.
\end{enumerate}
\end{nklogique}
